\documentclass[11pt]{article}
\usepackage[margin=1in]{geometry}
\usepackage{amsmath,amssymb,amsthm,booktabs,array,enumitem,microtype}
\usepackage[hidelinks]{hyperref}
\usepackage[T1]{fontenc}
\usepackage{lmodern}

\newtheorem{theorem}{Theorem}
\newtheorem{lemma}{Lemma}
\newcommand{\F}{\mathbb{F}}
\newcommand{\HGP}{\mathrm{HG}_{P}}

\title{The Proper Hat-Guessing Number of $K_5-e$}
\author{Matthew Protti\\\small Independent Researcher}
\date{2 September 2026\\\small Public research disclosure v0.1; independent review pending}

\begin{document}
\maketitle

\begin{abstract}
We prove that the proper hat-guessing number of the complete graph on five
vertices with one edge removed is eight. The upper bound follows from the
general inequality $\HGP(G)\le |V(G)|+\chi(G)-1$. For the lower bound, we
identify the eight colors with $\F_2^3$ and give two explicit rules for the
nonadjacent twin vertices. The residual proper colorings and the attainable
local views of the three clique vertices form a bipartite graph with left
degree three and right degree at most three. Hall's theorem therefore
completes the strategy. The only finite residue in the conceptual proof is a
seven-row table, comprising 42 direct evaluations. Dependency-free code
reconstructs the strategy and verifies all 8,400 proper colorings.
\end{abstract}

\section{Introduction}
Adriaensen, Bentley, Bishnoi, Cames van Batenburg, Kreiger, van der Kuil,
Mandal, Ramachandran, and Tuite introduced the proper-coloring variant of the
hat-guessing game. For a graph $G$, the parameter $\HGP(G)$ is the largest
number $q$ for which the players can pre-agree on local guessing functions so
that every proper $q$-coloring of $G$ makes at least one player guess
correctly. Their general bound
\[
   \HGP(G)\le |V(G)|+\chi(G)-1
\]
and their finite computations leave $\HGP(K_5-e)$ in the interval $[7,8]$.
We determine the exact value.

Let the two nonadjacent vertices of $K_5-e$ be called the \emph{twins}, and
let the other three vertices form the clique. Our proof first assigns the
twin guesses using a small geometry on $\F_2^3$. A matching argument then
assigns the clique guesses consistently.

\begin{theorem}\label{thm:main}
Let $K_5-e$ denote the complete graph on five vertices with one edge removed.
Then
\[
   \HGP(K_5-e)=8.
\]
\end{theorem}

\section{A finite geometry on eight colors}
Identify the color set with $V=\F_2^3$, written additively. For each nonzero
normal vector $n\in V$, let
\[
  U_n=\{x\in V:n\cdot x=0\}.
\]
These are the seven two-dimensional linear subspaces of $V$. Choose
$\delta(U_n)\notin U_n$ as follows.

\begin{center}
\begin{tabular}{c@{\qquad}c|c@{\qquad}c}
\toprule
$n$ & binary $n$ & $\delta(U_n)$ & binary $\delta$\\
\midrule
1 & 001 & 5 & 101\\
2 & 010 & 7 & 111\\
3 & 011 & 6 & 110\\
4 & 100 & 4 & 100\\
5 & 101 & 3 & 011\\
6 & 110 & 2 & 010\\
7 & 111 & 1 & 001\\
\bottomrule
\end{tabular}
\end{center}
The displayed choices satisfy $n\cdot\delta(U_n)=1$.

For nonzero $w\in V$, put $D_w=V\setminus\{0,w\}$ and define
\[
  \phi_w(r)=w+r+\delta(\langle w,r\rangle),\qquad r\in D_w.
\]

\begin{lemma}\label{lem:phi}
For every nonzero $w\in V$, the map $\phi_w$ is a fixed-point-free
permutation of $D_w$.
\end{lemma}

\begin{proof}
The complete cycle decompositions are
\begin{center}
\begin{tabular}{c|l}
\toprule
$w$ & cycles of $\phi_w$ on $D_w$\\
\midrule
1 & $(2\ 7\ 4)(3\ 6\ 5)$\\
2 & $(1\ 7\ 6)(3\ 5\ 4)$\\
3 & $(1\ 6\ 4)(2\ 5\ 7)$\\
4 & $(1\ 2\ 3)(5\ 6\ 7)$\\
5 & $(1\ 3\ 7)(2\ 4\ 6)$\\
6 & $(1\ 5\ 2)(3\ 4\ 7)$\\
7 & $(1\ 4\ 5)(2\ 6\ 3)$\\
\bottomrule
\end{tabular}
\end{center}
Thus each map is a permutation and has no fixed point. The table consists of
42 direct evaluations from the displayed definition of $\phi_w$.
\end{proof}

\section{The twin rules}
Label the clique vertices $2,3,4$. For an ordered triple
$T=(t_2,t_3,t_4)$ of pairwise distinct clique colors, define
\[
  \sigma(T)=t_2+t_3+t_4,
  \qquad
  U(T)=\langle t_2+t_3,t_2+t_4\rangle.
\]
The four points $t_2,t_3,t_4,\sigma(T)$ form the affine plane with direction
subspace $U(T)$. The twins guess
\[
  \alpha(T)=\sigma(T),
  \qquad
  \beta(T)=\sigma(T)+\delta(U(T)).
\]
The first guess is not one of the clique colors, and the second lies outside
the affine plane because $\delta(U(T))\notin U(T)$. Hence both guesses are
legal.

For each fixed clique triple, the twin colors may independently be any of the
five remaining colors. The rule $\alpha$ covers one row of the resulting
$5\times5$ array and $\beta$ covers one column. The twins therefore cover
$5+5-1=9$ pairs and leave 16 residual pairs. Since there are
$(8)_3=336$ ordered clique triples, the number of residual proper colorings is
$336\cdot16=5,376$.

\section{Residual incidence and Hall completion}
Form a bipartite graph $B=(L,R;E)$. The left part $L$ is the set of 5,376
residual proper colorings. The right part $R$ is the set of attainable local
views at the three clique vertices. Join a residual coloring to the three
clique views it induces. Every left vertex has degree three.

\begin{lemma}\label{lem:right}
Every right vertex of $B$ has degree at most three.
\end{lemma}

\begin{proof}
Fix a clique vertex and one of its local views. Let $x,y$ be the colors on the
other two clique vertices and let $a,b$ be the twin colors. Translate all
colors by $x$, and write
\[
  w=y+x,\qquad A=a+x,\qquad B=b+x.
\]
Then $w\ne0$ and $A,B\in D_w$. If the unseen clique vertex has translated
color $R$, the proper candidate set is $D_w\setminus\{A,B\}$ when $A\ne B$
and $D_w\setminus\{A\}$ when $A=B$.

For a candidate extension, twin 0 is correct precisely when $R=A+w$, and twin
1 is correct precisely when $\phi_w(R)=B$.

Suppose first that $A\ne B$. There are four proper candidates. If none were
twin-covered, then both distinguished candidates
$A+w$ and $\phi_w^{-1}(B)$ would lie in $\{A,B\}$. Since translation by $w$
has no fixed point, $A+w=B$. Since $\phi_w$ has no fixed point,
$\phi_w^{-1}(B)=A$. Hence
\[
  \phi_w(A)=B=A+w,
\]
contrary to
$\phi_w(A)=A+w+\delta(\langle w,A\rangle)$ and
$\delta(\langle w,A\rangle)\ne0$. At least one candidate is therefore
covered, leaving at most three residual extensions.

Now suppose $A=B$. There are five proper candidates. The candidates
$A+w$ and $\phi_w^{-1}(A)$ are allowed because translation by $w$ and
$\phi_w$ have no fixed points. They are distinct: equality would imply
$\phi_w(A+w)=A$, whereas
\[
  \phi_w(A+w)=A+\delta(\langle w,A\rangle)\ne A.
\]
Thus at least two candidates are covered, again leaving at most three
residual extensions.
\end{proof}

\begin{proof}[Proof of Theorem~\ref{thm:main}]
For every $S\subseteq L$, exactly $3|S|$ edges leave $S$. By
Lemma~\ref{lem:right}, every vertex of $N(S)$ is incident with at most three
of these edges. Thus
\[
  3|S|\le 3|N(S)|,
\]
so Hall's condition holds. Let $M$ be a matching saturating $L$. For each
matched residual coloring, assign the matched clique vertex to guess its own
color on that matched local view. A local view is matched at most once, so
these assignments are consistent; unmatched attainable views may receive any
legal guess. Every residual coloring is covered by a matched clique vertex,
and every nonresidual coloring was already covered by a twin. This is a
winning strategy with eight colors, so $\HGP(K_5-e)\ge8$.

The graph has five vertices and chromatic number four. The general upper
bound gives
\[
  \HGP(K_5-e)\le5+4-1=8.
\]
Therefore $\HGP(K_5-e)=8$.
\end{proof}

\section{Verification and scope}
The accompanying dependency-free Python code reconstructs the seven maps,
the residual incidence graph, a saturating matching, and a complete strategy.
It checks all 8,400 proper eight-colorings directly. These computations are
not premises of the conceptual proof.

This note settles only $K_5-e$. It does not solve the full $K_n-e$ family or
determine $\HGP(C_5)$.

\section*{AI-use disclosure}
OpenAI GPT-5.6 Pro was used substantially in literature search, construction
search, proof development, code generation, verification, adversarial
critique, and manuscript preparation. The human author selected and framed
the problem, directed the research, evaluated the outputs, required exact
checks, approved the final scope, and assumes responsibility for the work.

\begin{thebibliography}{9}
\bibitem{ABBCetal2026}
S. Adriaensen, P. Bentley, A. Bishnoi, W. Cames van Batenburg, M. Kreiger,
L. van der Kuil, S. Mandal, A. Ramachandran, and J. Tuite,
\emph{Hat guessing with proper colorings}, arXiv:2603.04909v4, 2026.
\end{thebibliography}

\end{document}
